\section{Models}
\label{models}

The study centres on three models that look at the problem in different ways. Logistic
regression on TF-IDF features is the standard supervised text baseline. The one-class SVM asks
whether fraud can be found without any fraud labels, by learning only what real ads look like. The
fine-tuned DistilBERT asks whether a pretrained language model reads ads better than word counts
do. Because DistilBERT and logistic regression turn out to make different mistakes, a fourth model
stacks the two. Every model goes through the protocol of Section~\ref{protocol} on the same five
splits, and a majority dummy fixes the floor: it predicts ``real'' for every ad, so its PR-AUC
equals the fraud rate (0.049).

\subsection{Logistic regression on TF-IDF features}

Logistic regression is the standard baseline for text classification. It gives every column of
the feature matrix (Table~\ref{tab:features}) its own weight, trains in about a minute and can be
read directly through its coefficients, which Section~\ref{explain} uses. We set
\texttt{class\_weight="balanced"}, which at 4.8\% positives makes a missed fraud cost about 20
times a false alarm during training. The tuning grid covers the regularisation strength
$C \in \{0.1, 1, 10\}$, a vocabulary of 10,000 or 30,000 terms, and single words or single words
plus word pairs.

\subsection{One-class SVM}
\label{ocsvm}

The one-class SVM~\cite{scholkopf2001} treats fraud as an anomaly. It is fitted on the real
training ads only and learns a boundary around the region they occupy; the fraud score of an ad is
how far it lies outside that boundary. The fraud labels are used only to leave the frauds out of
its training data and to choose its settings, so the model itself never sees a fraud example.

Two design choices decide whether this is a fair test, and our first version got both wrong. The
first is the input. The SVM's RBF kernel compares ads by their distance, and in the full feature
matrix the four numeric columns (ad and title length, exclamation marks, share of capitals) make up
45\% of each row's length, against under 5\% for the 20,000 text columns. Fed the full matrix, the
SVM mostly compared ad lengths. It therefore receives the TF-IDF text only. Because the RBF kernel
also needs short, dense input, the text is reduced to 100 dimensions with PCA (whose
\texttt{arpack} solver works on the sparse matrix directly), and each ad's vector is scaled to
unit length.

The second choice is the kernel width $\gamma$, which our first version left at its default. With
unit-length vectors, a small $\gamma$ makes the kernel broad, so the SVM judges an ad against the
bulk of all real ads, while a large $\gamma$ makes it local, so it judges an ad by whether it has
close neighbours among the real ads. A cross-validation map on the seed-0 training rows showed
that $\gamma$ decides almost everything: broad kernels ($\gamma \le 1$) rank frauds as \emph{more}
typical than real ads (ROC-AUC 0.23 to 0.30), while the score peaks around $\gamma = 30$ and falls
again beyond it. The other parameter, $\nu$ (roughly the share of training ads allowed outside
the boundary), barely matters. The final search therefore covers the whole grid
$\gamma \in \{3, 10, 30, 100, 300\}$, $\nu \in \{0.05, 0.2\}$ rather than a random sample of it.

\subsection{Fine-tuned DistilBERT}

DistilBERT~\cite{sanh2019} is a smaller and faster distillation of BERT~\cite{devlin2019}, a
transformer pretrained on large amounts of English text. Instead of counting words it reads each
ad in context, and fine-tuning updates all its weights on our training rows. It reads the first
256 tokens of each ad. Training uses AdamW with learning rate $2 \cdot 10^{-5}$ and weight decay
0.01, a linear warm-up over the first 10\% of steps followed by linear decay, batch size 128 and up
to 5 epochs. The loss uses the square root of the balanced class weights, a softer push towards
fraud than logistic regression gets, and the validation threshold sets the final trade-off. After
each epoch the validation PR-AUC is computed, the weights of the best epoch are kept, and training
stops after 2 epochs without improvement. The seed fixes the initialisation of the classification
head and the order of the training data, so every run can be repeated. DistilBERT runs on a GPU
server (notebook \texttt{03}) on the same saved splits.

\subsection{Stacking DistilBERT with logistic regression}
\label{stacking}

Section~\ref{explain} shows that DistilBERT and TF-IDF logistic regression catch different frauds.
When two models make different mistakes, combining them can beat both, so we stack
them~\cite{wolpert1992}: a small logistic regression whose only two inputs are the logits of the
two models' saved scores. The stacker has to be fitted on rows that the two models were not
trained on, because their scores on their own training rows are over-confident, so it is fitted
on the validation rows. Its threshold, which also has to be chosen on validation, comes from 5-fold
grouped out-of-fold predictions inside the validation rows, so it never sees a row the stacker was
fitted on. The test rows are scored once, as for every other model.

\subsection{Other models tried}

Five further models were run through the same protocol to check the choices above, and are
reported only briefly. BM25 logistic regression replaces TF-IDF by BM25~\cite{robertson2009}, the
weighting of classic search engines, in which repeated words have diminishing returns and long ads
are discounted. An MLP with one hidden layer and Complement Naive Bayes~\cite{rennie2003} use the
same features as logistic regression. Frozen sentence embeddings
(\texttt{all-MiniLM-L6-v2}~\cite{reimers2019,wang2020minilm}) with a logistic-regression head test the
cheapest use of a pretrained model, without fine-tuning. Finally, a zero-shot language model,
Qwen2.5-1.5B-Instruct~\cite{qwen2024}, is asked for each ad whether it is fraudulent, without any
training; its score is the probability of answering ``Yes''. All classical models use
scikit-learn~\cite{pedregosa2011}.
